\section{Finite cuts with compatible ends}\label{sec:finite-cuts}

The stretching construction supplies separated sets and a scalar-curvature
reserve. We now describe the second input: a cutting procedure that retains
both sides of every cut and controls how later cuts meet them. Its output is
a family of smooth boundaries and complete manifolds obtained by attaching
cylinders to the closed sides. The definitions below make its hypotheses
and output explicit; Appendices~\ref{sec:components}--\ref{sec:regularization}
give the proofs. Section~\ref{sec:countable} will separately justify its use
with countably many labels.

Let $X=M$ denote the original complete manifold, let $g$ be its original
metric, and write $d_0=d_g$ for its distance.  The metric used in an analytic
component can change; the map to $X$ always records its position in the
original manifold.  Fix finitely many of the supplied ball pairs
\[
 A_i=\overline B_{d_0}(x_i,D),\qquad
 B_i=X\setminus B_{d_0}(x_i,2D),\qquad 1\le i\le\ell,
\]
with $A_i,B_i$ nonempty. There is one label for each ball pair. An additional first label is radial.
Thus $N_*=\ell+1$, and the order is
radial, then $1,\ldots,\ell$.  Properness of $(X,g)$ makes each closed
outer ball compact.

The analytic input at stage zero is a complete smooth metric $h\ge g$ on $X$
and finitely many smooth positive circle warping functions for which the stabilized
scalar curvature is at least $.4$.  Each prescribed distance
$\dist_h(A_i,B_i)$ can be required to exceed any given finite number.
\cref{rounds:wide} supplies these data.  For finitely many requested
pairs its construction uses only finitely many compactly supported
increments and hence finitely many circle factors.  The countable
version has a locally finite list of nonunit warping functions and is used in
\cref{countable:construction}; the geometric statements below apply to
each finite collection of the relevant factors.

For a single smooth cut, let $u$ be a smooth function on an
$m$-manifold $E$ without boundary, with $0$ a regular value, and put
$L=u^{-1}(0)$.
Attach $L\times[0,\infty)$ to each closed side $\{\pm u\ge0\}$
along its boundary, and also retain $L$ as a separate space.
The two extended sides have dimension $m$, while the boundary has dimension
$m-1$. The cylinders replace the new boundaries by ends.
The following definition records these spaces and their incidences
after repeated cuts.

\begin{definition}[Component diagram]\label{component:category}
At stage $r$, an index is $\alpha\in\{+,-,*\}^{r}$.
Its entries record the first $r$ processed cuts.
A sign $+$ or $-$ selects the corresponding retained side of a cut,
and a star $*$ selects the zero set.
Indices are ordered by replacing signs with stars: an index obtained
from $\alpha$ by replacing $d\ge1$ signs by stars lies below $\alpha$.
A stage-$r$ component diagram consists of the following data, with the
following properties.
\begin{enumerate}[label=(\roman*),leftmargin=*]
\item \emph{Spaces.} For each index $\alpha$, the space $E_\alpha$ is a
smooth second countable manifold without boundary,
possibly empty or disconnected; when nonempty its base dimension is
$m(\alpha)=n-\#\{i:\alpha_i=*\}$.  We consider its connected components.
\item \emph{Product tails.}
If $B$ is a component of an index $\beta$ and $\alpha$ is obtained from $\beta$
by replacing $d$ stars with signs, then in $E_\alpha$ the component $B$ has
exactly one incident component $E$ and
one corresponding occurrence of a full product tail
\begin{equation}\label{component:tail}
 \iota_{BE}:B\times(T_1,\infty)\times\cdots\times(T_d,\infty)
       \longrightarrow E.
\end{equation}
\(B\) is the base of this product end of \(E\), with \(d\) normal factors.
\(E\) is the corresponding incident higher-dimensional component.
For \(d=1\), the product end has one normal coordinate.
The normal coordinates are labelled by the changed index entries.
Thresholds may depend on the incidence.
\item \emph{Intersections and commutation.} For a single entry, tails
from distinct end base components are disjoint.  Intersections of sufficiently
far tails for different entries are exactly the tails over their common
further end bases.  The identifications commute: composing end base maps in
any order gives the same occurrence and the same individually labelled
normal coordinates.  After slightly raising the entry thresholds,
each closed tail is properly embedded.
\item \emph{Position maps.} Each component $E$ is equipped with a continuous
map $q_E:E\to X$.
\item \emph{After the radial stage}, in addition:
every component has only finitely many incident end bases and incident higher-dimensional components, including
through chains; the complement of its sufficiently far product tails with one normal factor
is compact.  On each sufficiently far end base product the maps $q_E$ satisfy
$q_E(\iota_{BE}(b,\mathbf t))=q_B(b)$.
\end{enumerate}
A metric, scalar function, warping function, or set is \emph{compatible} if, in every
sufficiently far tail, it is exactly the product or pullback of the
corresponding end base datum.  For metrics the normal factors are
$\dd t_1^2+\cdots+\dd t_d^2$; all warping functions are independent of the normal
coordinates there.  The thresholds for different data may differ.
\end{definition}

There is no boundary in an analytic component.  Its end base is represented by
an entire cylindrical end, not by a boundary subset.  Finite boundaries
appear only when we truncate these ends.

A \emph{split along regular zero sets} applies the preceding construction using
compatible smooth defining functions whose zero sets are regular on
every index space. Each zero set gives the space indexed by a new star.
In inherited ends the operation is the corresponding end base operation
times the old normal factors. The proof that these operations preserve
the component diagram is \cref{component:split}.

For the diagrams constructed here, these maps are defined recursively,
starting with $q_X=\Id_X$.  In the split of \cref{component:split}, let
$A$ be a new stage-$(r+1)$ component with stage-$r$ parent $E$.
Define $p_A:A\to E$ by inclusion on a zero set or the bulk of a joined
side, and by $p_A(x,t)=x$ on each newly attached cylinder.
Set $q_A=q_E\circ p_A$.  Thus $q_A$ is pulled back from the base on
each fresh cylinder.  In a sufficiently far inherited tail of $A$,
let $C$ be the corresponding new end base with old parent $B$, and let
$I$ denote the old normal intervals.  Under the identifications of
these tails with $C\times I$ and $B\times I$, the map $p_A$ is
$p_C\times\Id_I$; explicitly,
\[
 p_A\circ\iota_{CA}=\iota_{BE}\circ(p_C\times\Id_I),\qquad
 q_A\bigl(\iota_{CA}(c,\mathbf t)\bigr)=q_C(c).
\]
Thus the recursion preserves compatibility.  The proper maps $q_r$
on finite reconstructions are constructed in
\cref{component:reconstruction}.

Figure~\ref{fig:regular-component-split} depicts one extended side and the
separate zero set.
\begin{figure}[!b]
\centering
\begin{tikzpicture}[x=1cm,y=1cm,>=Stealth,font=\small]
  \fill[blue!7] (0,0) rectangle (2.3,2.2);
  \fill[orange!12] (2.3,0) rectangle (5,2.2);
  \draw[thick] (0,0) -- (5,0);
  \draw[thick] (0,2.2) -- (5,2.2);
  \draw[densely dashed,thick] (2.3,0) -- (2.3,2.2);
  \node[align=center] at (1.1,1.2) {retained\\closed side};
  \node[align=center] at (3.6,1.2) {$L\times[0,\infty)$\\attached\\cylinder};
  \draw[->,thick] (4.4,.5) -- (5.3,.5) node[right] {$s$};
  \node[align=center,below] at (2.3,-.1) {internal attachment interface};
  \draw[very thick,teal!70!black] (7,0) -- (7,2.2);
  \node[align=center,right] at (7.2,1.1) {separate zero set $L$\\base dimension $m-1$};
  \node[above] at (2.5,2.7) {either sign side: base dimension $m$};
\end{tikzpicture}
\caption{A regular cut retains each closed sign side with a full cylinder
attached, and also retains the boundary as a separate zero set of one lower base
dimension. The attachment interface is internal to the extended side. In old
product ends this entire operation is a product. The raw metric and its
later smoothing are defined in Appendices~\ref{sec:cut} and~\ref{sec:regularization}.}
\label{fig:regular-component-split}
\end{figure}


The label set and the index depth are finite in this section. The family
of connected components may be countable; finiteness is required only for
the directed chains through end bases and higher-dimensional components
incident to each component. The
downward diagram of a component consists of that component and all its further
end bases. Numerical
requests are allocated downward through these chains before actual data
are constructed upward from end bases.

We now specify the analytic data that must survive each split.
For the finite-label statements in this section and the appendices we
use the following constants:
\begin{equation}\label{component:constants}
\begin{gathered}
 \lambda_r=\frac14-\frac{r}{8N_*},\qquad
 b_r=1.12-\frac{.05r}{N_*},\qquad
 a_r=.60+\frac{.03r}{N_*}\quad(0\le r\le N_*),\\
 \pi_0=\frac{1}{10N_*},\qquad k=399,\qquad
 C=\frac{4k}{k+1}=3.99,\qquad
 e_0=\frac{\lambda_r-\lambda_{r+1}}{100}
     =\frac{1}{800N_*}\quad(r<N_*).
\end{gathered}
\end{equation}
Here $k$ is the number of circle factors added at each split; the initial
number may be different. The choice $k=399$ realizes the coefficient
$C=4k/(k+1)$ in the scalar-curvature identity of \Cref{regular:step}.
With $C=3.99$, the boundary and retained-side estimates leave gradient
coefficients $C-2=1.99$ and $C-3.05=.94$, respectively
(\Cref{cut:raw-domination}); these supply the reserve for enlargement of the prescribed open sets.
In particular $\lambda_r\ge1/8$, $a_r\le.63$,
and $1.07\le b_r\le1.12$.  The common floor $\pi_0$ and the equally
spaced stage parameters are bookkeeping for a finite list, not uniform
lower bounds for the countable construction. Section~\ref{sec:countable}
and Appendix~\ref{sec:regularization} write the same constants in terms
of increasing budget parameters \(t_r\in[0,1]\), with increments
\(\Delta_r=t_{r+1}-t_r\), in \eqref{countable:constants}. The equally
spaced choice \(t_r=r/N_*\) gives back \eqref{component:constants}, with
\(\Delta_r=1/N_*\) and \(e_0=\Delta_r/800\), and gives the number
\(\eta_r=\Delta_r/200=1/(200N_*)\) that appears in \Cref{finite:cut}.

\begin{definition}[Analytic invariant]\label{component:invariant}
At stage $r$, write $J_r$ for the unprocessed labels.  On a component $E$
of base dimension $m$, retain its base metric and all circle warping functions.
Write $R_E$ for the scalar curvature of that full warped metric and
$\omega_E$ for the product of its warping functions, suppressing fixed
circle-volume constants.  Gradients of invariant functions are base
gradients; their integrals use $\omega_E\,\dd\vol_E$, and their
Laplacian is the corresponding weighted base Laplacian.
For every $j\in J_r$ the data comprise a compatible continuous function
$\mu_j:E\to[-\infty,\infty]$ whose finite truncations
$\max\{-c,\min\{\mu_j,c\}\}$, $0<c<\infty$, are locally Lipschitz,
a compatible smooth potential $\Pi_j\ge\pi_0$, and compatible prescribed
open sets $U_j^+,U_j^-$ satisfying
\begin{equation}\label{component:open-set-poles}
 \overline{U_j^\pm}\subset
 \operatorname{int}_E\{\mu_j=\pm\infty\}.
\end{equation}
On the open finite-value locus $\{x\in E:|\mu_j(x)|<\infty\}$ the function
obeys, almost everywhere,
\begin{equation}\label{component:function}
 2|\grad\mu_j|\le a_r\mu_j^2+\Pi_j.
\end{equation}
Set
\begin{equation}\label{component:required-set}
 G_j=E\setminus(U_j^+\cup U_j^-),\qquad
 D_E=\bigl\{x:\#\{j\in J_r:x\in G_j\}\ge m-1\bigr\}.
\end{equation}
The required scalar bound is
\begin{equation}\label{component:curvature}
 R_E\ge\lambda_r+b_r\sum_{j\in J_r}\Pi_j
       \quad\hbox{on }D_E.
\end{equation}
For an unprocessed label $i$, the prescribed open sets contain
$q_E^{-1}(A_i)\subset U_i^+$ and
$q_E^{-1}(B_i)\subset U_i^-$.  Enlarging prescribed open sets must preserve these
conditions.  The sets $D_E$ are closed: they are finite unions of finite intersections
of the closed sets $G_j$ (and equal $E$ when the count threshold is
nonpositive).  All data are exact products
in the sense of \cref{component:category}; at stage zero the single
space is $X$, with no product-end requirement.
\end{definition}

The threshold $m-1$ is preserved by the two kinds of output of a
split. A side keeps its dimension and loses one unprocessed label;
a zero set loses both one dimension and that label. The set $H_E^*$ where at least $m-2$ remaining labels are undecided contains the projections of all new required
sets, even after the remaining prescribed open sets have been enlarged.

\begin{lemma}[The two count inclusions]\label{component:counts}
Use the required sets and compatible prescribed open sets of
\Cref{component:invariant}. If $B$ is the base of a product end of $E$ with $d$ normal factors,
the specified tail projection in \eqref{component:tail} sends
$D_E$ in that tail into $D_B$.
At a split of label $j_0$ on a parent of dimension $m$, suppose the
remaining prescribed open sets have been compatibly enlarged, giving $G'_j\subset G_j$ for
$j\ne j_0$, and define the closed set
\begin{equation}\label{component:count-set}
 H_E^*=\bigl\{x:\#\{j\in J_r\setminus\{j_0\}:x\in G'_j\}
                      \ge m-2\bigr\}.
\end{equation}
Suppose a regular zero set $L$ lies in $\operatorname{int}G_{j_0}$.
Pull the remaining prescribed open sets to a new side or component of the zero set $A$ by bulk
inclusion, boundary restriction, and projection of the attached cylinder.
Define $D_A$ using these prescribed open sets, the remaining labels, and the dimension
of $A$, as in \eqref{component:required-set}.
Then the bulk inclusion, boundary restriction, or cylinder projection
that defines $D_A$ sends it into both $H_E^*$ and the parent's old
set $D_E$.
The sets $H_E^*$ likewise project into $H_B^*$ on further end bases $B$.
\end{lemma}

\begin{proof}
In an exact product tail all label memberships pull back unchanged,
whereas the end base dimension is $m-d$.  A count at least $m-1$ is at
least $(m-d)-1$, proving the first assertion.  The same argument
with thresholds $m-2$ and $(m-d)-2$ proves the set $H_E^*$ assertion.
On a new side, whose dimension is $m$, a point of $D_A$ has at least
$m-1$ remaining undecided labels.  This puts its projection in the
set $H_E^*$ and, because $G'_j\subset G_j$, in the old required set.
On a component of the zero set, whose dimension is $m-1$, the required count is
$m-2$.  Its projection lies on $L\subset G_{j_0}$, so the processed
label contributes one further old undecided label.  The old count
is therefore at least $m-1$.  The cylinder projection has exactly
the same remaining memberships as its base, so the side argument
also covers all cylinder points.  These are pointwise count
statements.
\end{proof}

\begin{lemma}[Initialization and radial confinement]\label{component:initialization}
For the stage-zero metric, warping functions, and ball pairs fixed above, assume
that the stabilized scalar curvature is at least $.4$ and that each
label pair satisfies
\[
 \dist_h(A_i,B_i)>T_0+4,\qquad
 T_0=2\pi\sqrt{2/\pi_0}.
\]
Then the given metric and warping functions admit the functions, prescribed open sets, and
potentials of \cref{component:invariant} at stage zero with
$\Pi_j=\pi_0$.  On a noncompact $X$ the radial prescribed open sets can be chosen
as separating inverse images of intervals of a smooth proper function,
with alternating signs,
with compact smooth barrier domains around their closures lying in
the same pole plateaus.  Every regular radial partition respecting
these signs has locally finite compact boundary components and compact
closed-side components before collar attachment.  On a compact $X$
a decided radial label with empty boundary has the same properties.
\end{lemma}


\begin{theorem}[Finite compatible cutting]\label{finite:cut}
Fix $N_*$ and the stage constants in \eqref{component:constants}. Let
$0\le r<N_*$ and let a stage-$r$ diagram arise from the regular
splits of \Cref{component:category}, with compatible smooth base metrics
and a finite common list of positive circle warping functions satisfying
\Cref{component:invariant}. After the radial stage, assume that each
component is complete, has compact remainder outside its finitely many
product ends, and has the exact metric and warping function identifications in
\Cref{component:category}. Let $j_0$ be the next label, and assume that
every nonempty parent has base dimension at least two.

The diagram admits a compatible split along regular zero sets of $j_0$ with the
following properties.
\begin{enumerate}[label=(\roman*),leftmargin=*]
\item On each parent $E$, the smooth defining function has the
prescribed signs on the current closures of the prescribed open sets, and its regular
zero set lies in $\operatorname{int}G_{j_0}$. The new spaces are
that boundary and the two closed sides with their full boundary cylinders
attached. Their maps to $X$ are the continuous maps $q_A=q_E\circ p_A$
specified above, including the commuting inherited tail projections.
\item The remaining parent functions and potentials may change, with
prescribed open sets only enlarged and potentials only increased. Their restrictions
and cylinder pullbacks admit compatible smooth metrics and potentials
and exactly $399$ new positive circle factors satisfying the full
stage-$(r+1)$ invariant. Earlier forced signs persist. The new components
are complete with compact remainders and exact product ends. This
includes a possible one-dimensional output, whose scalar bound holds
everywhere.
\item Every old warping function identically equal to $1$ stays fixed. A parent's
downward source diagram consists of that component and all its further
end bases, together with every specified product occurrence
\eqref{component:tail} whose end base and incident higher-dimensional component both belong to this diagram.
An end base copy means the data on the end base pulled back to such an
occurrence. If a remaining label has the same constant pole function,
full corresponding prescribed open set, and spatially constant potential $c_j$ on
every component of this finite diagram, and hence on all their end base copies,
its pole and prescribed open set can stay fixed
and its new potential can equal $(1+3\eta_r/2)c_j$ exactly, where
$\eta_r=1/(200N_*)$.
\end{enumerate}
At stage zero the parent may instead be an arbitrary complete smooth
manifold. With the same scalar and separation assumptions and the
radial initialization of \Cref{component:initialization}, the first split
has all the stated conclusions. Its boundary and closed-side components
are compact before cylinder attachment and form locally finite families.
\end{theorem}

\begin{proof}
Realize the processed function by a smooth finite prescribed weighted mean curvature and take a
local minimizer of weighted perimeter minus the weighted volume
integral of that function, as in \cref{cut:prescribed-curvature,cut:bv-compactness}. The Caccioppoli set has the
forced signs of \cref{cut:barriers}; the support of its perimeter measure
has the regular-part geometry of \cref{cut:gmt} but may still have
singularities. After the
radial cut, the generic prescribed weighted mean curvature choice in
\cref{cut:generic-uniqueness} makes these boundary supports converge to the
corresponding end base products at component ends
(\cref{cut:end-convergence}). On the initial radial parent use shift
zero and a Caccioppoli set from compactness; compact barriers replace this
post-radial growth and uniqueness argument.

The boundary, and each retained side with its attached cylinder,
carry the quadratic-form lower bounds of \cref{cut:raw-domination}.
These bounds include a positive gradient reserve and, on a joined
side, the scalar contribution of its internal attachment interface.
They are not yet pointwise scalar bounds on smooth new components.
Keep these boundary supports and their prescribed weighted mean curvatures fixed. The constructions of
\Cref{local:smooth-approximations,local:sharp,enlargement:global} change only the remaining
functions, enlarge their prescribed open sets, and add nonnegative potentials. Their
form cost is paid by the gradient reserve and the gap
$\lambda_r-\lambda_{r+1}$, and the updated set $H_E^*$ of
\eqref{component:count-set} misses the singular set. Thus the future required
regions can be controlled through regular portions of the boundary support.
When the parent dimension is at most seven the
singular set is empty, and no singular forcing is needed; compactness on
the resulting smooth regions is the smooth case of
\cref{enlargement:compactness}. More generally this case with a smooth limiting boundary applies
whenever the limiting boundary is smooth, including products of smooth end bases
with Euclidean factors even when their total dimension is at least eight.
At the first radial split the same arguments
are applied on the locally finite compact transition blocks, as specified
in \cref{enlargement:global}.

The approximation near $H_E^*$ in \cref{regular:partition} now replaces
the boundary support by a smooth regular boundary. Its defining functions are
exact products in the old ends. The comparison with the boundary support on
the set $H_E^*$, followed by the interface smoothing of \cref{regular:interface},
transfers the form bound to the smooth split. The positive
supersolutions of \cref{regular:correction,regular:solutions}
then supply $399$ compatible new circle factors; the scalar identity
in \cref{regular:step} converts that form bound into the pointwise
stage-$(r+1)$ scalar inequality. Throughout this transfer the parent
scalar inequality is used only over $D_E$: \cref{component:counts}
places every new required region inside that old region as well as
inside the set $H_E^*$.

\Cref{component:split,component:geometry} give the
new geometry; \cref{component:initialization} gives the radial confinement.
Unit entries are fixed by restriction, product extension, and interface
smoothing. For a remaining label with the same constant pole function, full
corresponding prescribed open set, and constant potential throughout its
source diagram,
we omit its forcing updates, and
\cref{regular:constant-preservation,regular:functions} reproduce its chosen
constant potential exactly. The choice of $\eta_r$ satisfies the strict
stage margins in \eqref{regular:eta}.
\end{proof}
